\section{Background and Related Work}
\label{sec:background}

\subsection{Majority-Inverter Graphs}
\label{sec:background:mig}

A Majority-Inverter Graph (MIG)~\cite{amaru2014mig} is a logic
representation in which every internal node is a 3-input majority
gate $\maj3(a, b, c) = ab \lor bc \lor ac$, with optional
complementation on each input edge. MIG synthesis flows such as the
DAC'19 compatible pipeline of~\cite{amaru2019synth} apply algebraic
rewriting, resubstitution, and refactoring passes that exploit
majority-specific identities to reduce node count and depth.

\subsection{Differential fault analysis and the role of representation}
\label{sec:background:dfa}

Differential Fault Analysis (DFA) was introduced
by Biham and Shamir~\cite{biham1997dfa} as a key-recovery attack
against block ciphers under fault injection. The attacker induces a
fault during an encryption, observes the resulting incorrect output,
and uses the difference between correct and faulty ciphertexts to
constrain the secret key. The seminal AES DFA
of~\cite{piret2003dfa} requires as few as two faulty ciphertexts
under a single-byte fault at the right intermediate position; later
work has extended DFA to most major symmetric ciphers and to a
variety of fault models~\cite{bareltal2006fault, joye2012faultattacks}.

The exploitability of a given fault, and therefore the cost of an
end-to-end DFA attack, depends on three quantitative variables
that are functions not only of the cipher but of its hardware
implementation: the number of distinct fault sites accessible to
the adversary, the per-vector probability that an injected fault
produces a detectable output deviation, and the algebraic structure
of the resulting differences. The first two of these vary with the
synthesis flow. We measure them; we do not address the third.

\subsection{Fault analysis tools and prior characterization}
\label{sec:background:tools}

Pre-silicon fault analysis frameworks evaluate the behavior of
synthesized netlists against a parametrized adversary. SYNFI of
Nasahl et al.~\cite{nasahl2022synfi} performs formal pre-silicon
fault analysis on industrial netlists, focusing on whether
designer-inserted countermeasures survive logic-synthesis
optimization. SCFI of~\cite{nasahl2022scfi} hardens
finite-state-machine control flow against multi-fault attacks.
FIVER~\cite{richter2021fiver} and related tools target similar
goals from different angles.

These tools all start with a single, fixed netlist and characterize
how an attacker can exploit it. Our study is complementary: we hold
the cipher and the FPGA target fixed and vary the synthesis flow,
measuring how the choice of gate-level representation moves the
fault observability metrics that those tools take as input.
We do not address countermeasure design or formal fault verification.

\subsection{Higher-arity majority and AQFP}
\label{sec:background:majn}

The mockturtle library~\cite{soeken2018mockturtle} provides a
\texttt{majority5} generator that decomposes MAJ-5 into four MAJ-3
gates following the construction of~\cite{amaru2014mig}. The
\texttt{aqfp\_network} variant supports native n-ary majority nodes
for Adiabatic Quantum-Flux Parametron (AQFP) targets, where
higher-arity majority is a natural primitive of the cell
library~\cite{takeuchi2014aqfp}. No production
cipher-synthesis flow currently targets MAJ-$n$ as a primitive;
\S\ref{sec:results:majn} discusses why this matters for fault
observability.
